\documentclass[amssymb,pra,12pt,aps,notitlepage]{revtex4-1}
\usepackage{mathptmx,amsmath, multirow, graphicx}
\usepackage{parskip}
\begin{document}

\title{Putnam POTD \today}
\maketitle

\section{Problem}
(1990 A3)Prove that any convex pentagon whose vertices (no three of
which are collinear) have integer coordinates must have area greater than
or equal to 5/2.

\section{Solution}

Recall from Pick's theorem that for a convex polygon $P$, 
Let $I$ be the set of grid points inside $P$, and $B$ be the 
set of grid points on the boundary of $P$, then 
\[Area(P) = \left|I\right| + \frac{1}{2} \left|B\right| - 1.\]

In our case, $B = 5$, thus we only need to show that 
$I + \frac{3}{2} \ge \frac{5}{2}$, so 
we need to show $I \ge 1$. 

Since all coordinates are integers, they can have four possible parity \\
configurations:
$(E,E),(O,O),(O,E),(E,O)$, which can be seen as 
4 pigeonholes, and since there are 
 $5$ points, so there must be two points with the same configuration. 
For two points with the same configuration, their midpoint must be a grid point. 

There are two cases: If the two points are diagonally across from each other,
their grid point is in the interior, 
then there is a grid point inside and we are done. \\Otherwise, the two points 
are on the boundary, and their midpoint is also on the boundary,
which reduces the problem to a smaller subproblem, as 
we have another convex pentagon on grid points again. 

From a high level perspective, this process will terminante eventually,
hence there will always be an interior point, and thus 
$Area(P) \ge \frac{5}{2}$. In fact, the construction for 
$Area(P) = \frac{5}{2}$ is given by 
the polygon $(0,0),(1,0),(0,1),(1,1),(1,2)$. 



\end{document}